\section{Background}
\label{sec:bg}

\newcommand{\good}[1]{\leavevmode\color{Green}{#1}}
\begin{table}
\caption{Comparison of Decentralized Storage Networks}
\label{tab:dsn}
% \begin{tabularx}{\linewidth}{XXXXX}
%     \toprule
%     Works       & Systematic placement  & Configurable redundancy   & Friendly to light nodes   & Effective redundancy  \\
%     \midrule
%     Filecoin    & No                    & \good{Yes}                & No                        & \good{Low}            \\
%     Arweave     & No                    & No                        & No                        & High                  \\
%     Storj       & Centralized           & \good{Yes}                & \good{Yes}                & \good{Low}            \\
%     Swarm       & \good{Decentralized}  & \good{Yes}                & \good{Yes}                & High                  \\
%     \sys        & \good{Decentralized}  & \good{Yes}                & \good{Yes}                & \good{Low}            \\
%     \bottomrule
% \end{tabularx}
\begin{tabularx}{\linewidth}{lXXX}
    \toprule
    Works       & Friendly to light nodes   & Systematic placement  & Effective redundancy  \\
    \midrule
    Filecoin    & No                        & No                    & \good{Low}            \\
    Arweave     & No                        & No                    & High                  \\
    Storj       & \good{Yes}                & Centralized           & \good{Low}            \\
    Swarm       & \good{Yes}                & \good{Decentralized}  & High                  \\
    \sys        & \good{Yes}                & \good{Decentralized}  & \good{Low}            \\
    \bottomrule
\end{tabularx}
\end{table}

\subsection{Decentralized Storage Networks}

Decentralized storage networks are widely used to outsource bulk data from blockchains.
\xxx{Facts.}

\xxx{Arch: read/update metadata on chain, read/append content data off chain.}

DSNs are under the same critical attack vector as the blockchains.
By utilizing DSNs, the durability of DSNs becomes part of the safety guarantee of blockchains themselves.
Proof of storage can rapidly detect faults and allows DSNs to repair (relatively) timely.

\paragraph{Threat model.}
The time is divided into \emph{proof periods}.
Storage providers submit proof of storage for the contents they stored in every proof period.
The proofs attest:
\begin{inparaenum}
    \item the stored contents are durable with sufficient redundancy, which enables external auditing, and
    \item the utilized capacities of storage providers.
\end{inparaenum}
The storage capacity must be correctly utilized in a proof period, even if it belongs to a Byzantine faulty storage provider, as long as a proof of storage is submitted for it.
On the other hand, any faulty storage capacity can only be detected at the end of the proof period because of the lack of proof of storage.
We assume that at most $\frac{1}{3}$ of the storage capacity can be faulty in a proof period.
We also assume that in a proof period the system is able to finish repairing for any repairable faults.
For a system that only employs replication as redundancy, a fault is repairable means not all copies of the data are lost.

\subsection{The Problem with DHT}

Some storage providers, especially the ones with smaller capacities, are full before the others.

\xxx{A graph that samples utilization of random storage providers at a time point, showing imbalance.}

\xxx{Or a graph that traces utilization of random storage providers over time, showing inconsistent growing speed.}

\xxx{Or a graph that where x-axis is the overall utilization and y-axis is the portion of nodes with at least X utilization, showing many nodes are saturated while system is still at a low utilization.}

Even if the storage providers are already full, DHT still ``can't help'' to push more data to them because who to push is strictly enforced by the DHT protocol and the distance algorithm.

\paragraph{Simply asking users to store ``somewhere else'' is not a solution.}
One simple workaround is to manipulate the content hash by \eg equipping the content with a nonce, and choose a hash where the closest storage providers in the DHT have sufficient remaining storage capacity.

\begin{figure}
    \centering
    \includegraphics[width=\linewidth]{data/simulation/graphs/ingest-reject}
    \caption{Number of attempts per ingestion at different utilization.}
    \label{fig:ingest-reject}
\end{figure}

As \autoref{fig:ingest-reject} shows, this does not work well for ingestion.
\xxx{Not a good graph, maybe change to utilization vs. average number of attempts.}

And this trick completely doesn't work when the DSN scales down.

\paragraph{Virtual node.}
Balance storage providers' capacity by splitting larger storage providers into smaller virtual ones.

As we will see in \autoref{sec:overview}, this only solves half of the problem.
Even in a DSN where most of the storage providers have equal capacity still suffers from low utilization.

Virtual nodes also introduce artificial correlation that DHT protocol is not aware and cannot account for.
DHT protocol could place redundancy of the same content repeatedly on the same physical storage provider.
\xxx{Some simulation if desired.}